% GENERATED FILE. Edit canonical lessons, then run npm run generate:books.
% canonical-source-hash: fnv1a64:6f5c46cac7e32c32
% canonical-lessons: LA-C315-naufragium, LA-C315-naufragus, LA-C315-mergo, LA-C315-obruo, LA-C315-calamitas

\chapter{Naufragium}
\label{ch:la-naufragium}
\hlchaptermodality{315}

\begin{chapteropening}
\textbf{By the end of this chapter:} \emph{I can tell the story of a shipwreck: the wreck, shipwrecked, to sink, to overwhelm and a disaster.}

Complete the last lesson of chapter 315: I can tell the story of a shipwreck: the wreck, shipwrecked, to sink, to overwhelm and a disaster.
\end{chapteropening}

\section[naufragium, naufragiī]{naufragium, naufragiī — a shipwreck}

\label{lesson:LA-C315-naufragium}



\begin{quote}
Before the new one: say the Latin for a blow, then the Latin for wounded.
\end{quote}

\subsection*{You'll want to know: naufragium, naufragiī}
\textbf{naufragium, naufragiī} — \textquotedblleft{}a shipwreck\textquotedblright{}.

\begin{quote}
Naufragium prope lītus factum est.
\end{quote}

\begin{quote}
\emph{A shipwreck happened near the shore.}
\end{quote}

Neuter, built from the words for 'ship' and 'to break'. \emph{Naufragium facere}: to be shipwrecked.

\begin{cousinweb}
\textbf{naufragium} stands for an older \emph{*nāvifragium}, from \emph{nāvis}, \textquotedblleft{}ship\textquotedblright{}, and \emph{frangō}, \textquotedblleft{}to break\textquotedblright{} (Lewis \& Short): a ship-breaking.
\end{cousinweb}

\begin{grammarlens}[title={What you've built}]
One more word for telling, explaining and describing at B1.
\end{grammarlens}

\subsection*{Your turn}
\begin{itemize}
\raggedright
  \item \emph{Say it:} \emph{naufragium, naufragiī}
  \item \emph{Say it:} \emph{naufragium, naufragiī}, once more
  \item \emph{Say it:} a sentence of your own with \emph{naufragium, naufragiī}
  \item \emph{Recall:} say the Latin for a blow, then the Latin for wounded, then say \emph{naufragium, naufragiī} again
\end{itemize}

\begin{tcolorbox}[breakable,colback=teal!4,colframe=teal!35!black,title={Before you move on}]
What does naufragium, naufragiī mean? (\textquotedblleft{}A shipwreck\textquotedblright{}.) Say it once more.
\end{tcolorbox}
\section[naufragus, naufraga, naufragum]{naufragus, naufraga, naufragum — shipwrecked}

\label{lesson:LA-C315-naufragus}



\begin{quote}
Before the new one: say the Latin for wounded, then the Latin for a shipwreck.
\end{quote}

\subsection*{You'll want to know: naufragus, naufraga, naufragum}
\textbf{naufragus, naufraga, naufragum} — \textquotedblleft{}shipwrecked\textquotedblright{}.

\begin{quote}
Nautae naufragī in lītore sedēbant.
\end{quote}

\begin{quote}
\emph{The shipwrecked sailors were sitting on the shore.}
\end{quote}

As a noun, \emph{naufragus} is a shipwrecked man; beggars told of their wrecks to win coins.

\begin{cousinweb}
\textbf{naufragus} joins \emph{nāvis}, \textquotedblleft{}ship\textquotedblright{}, and \emph{frangō}, \textquotedblleft{}to break\textquotedblright{} (EtymDB): one whose ship is broken.
\end{cousinweb}

\begin{grammarlens}[title={What you've built}]
One more word for telling, explaining and describing at B1.
\end{grammarlens}

\subsection*{Your turn}
\begin{itemize}
\raggedright
  \item \emph{Say it:} \emph{naufragus, naufraga, naufragum}
  \item \emph{Say it:} \emph{naufragus, naufraga, naufragum}, once more
  \item \emph{Say it:} a sentence of your own with \emph{naufragus, naufraga, naufragum}
  \item \emph{Recall:} say the Latin for wounded, then the Latin for a shipwreck, then say \emph{naufragus, naufraga, naufragum} again
\end{itemize}

\begin{tcolorbox}[breakable,colback=teal!4,colframe=teal!35!black,title={Before you move on}]
What does naufragus, naufraga, naufragum mean? (\textquotedblleft{}Shipwrecked\textquotedblright{}.) Say it once more.
\end{tcolorbox}
\section[mergō, mergere]{mergō, mergere — to sink, to plunge, to dip}

\label{lesson:LA-C315-mergo}



\begin{quote}
Before the new one: say the Latin for a shipwreck, then the Latin for shipwrecked.
\end{quote}

\subsection*{You'll want to know: mergō, mergere}
\textbf{mergō, mergere} — \textquotedblleft{}to sink, to plunge, to dip\textquotedblright{}.

\begin{quote}
Undae nāvem parvam mersērunt.
\end{quote}

\begin{quote}
\emph{The waves sank the little ship.}
\end{quote}

\emph{mergō, mergere, mersī, mersum}. English \emph{merge} and \emph{submerge} come from it.

\begin{cousinweb}
\textbf{mergō} is an inherited verb that Lewis \& Short connect with a Sanskrit verb meaning \textquotedblleft{}to dip\textquotedblright{}. English \emph{merge}, \emph{submerge} and \emph{immerse} come from it.
\end{cousinweb}

\begin{grammarlens}[title={What you've built}]
One more word for telling, explaining and describing at B1.
\end{grammarlens}

\subsection*{Your turn}
\begin{itemize}
\raggedright
  \item \emph{Say it:} \emph{mergō, mergere}
  \item \emph{Say it:} \emph{mergō, mergere}, once more
  \item \emph{Say it:} a sentence of your own with \emph{mergō, mergere}
  \item \emph{Recall:} say the Latin for a shipwreck, then the Latin for shipwrecked, then say \emph{mergō, mergere} again
\end{itemize}

\begin{tcolorbox}[breakable,colback=teal!4,colframe=teal!35!black,title={Before you move on}]
What does mergō, mergere mean? (\textquotedblleft{}To sink, to plunge, to dip\textquotedblright{}.) Say it once more.
\end{tcolorbox}
\section[obruō, obruere]{obruō, obruere — to overwhelm, to bury}

\label{lesson:LA-C315-obruo}



\begin{quote}
Before the new one: say the Latin for shipwrecked, then the Latin for to sink.
\end{quote}

\subsection*{You'll want to know: obruō, obruere}
\textbf{obruō, obruere} — \textquotedblleft{}to overwhelm, to bury\textquotedblright{}.

\begin{quote}
Tempestās nautās obruit.
\end{quote}

\begin{quote}
\emph{The storm overwhelmed the sailors.}
\end{quote}

\emph{obruō, obruere, obruī, obrutum}: \emph{ob} + \emph{ruō}, to rush over. Also of earth and ash: a town buried under ash.

\begin{cousinweb}
\textbf{obruō} is built on \emph{ruō}, \textquotedblleft{}to rush, to fall\textquotedblright{}, with \emph{ob-} (EtymDB): to fall upon and cover over.
\end{cousinweb}

\begin{grammarlens}[title={What you've built}]
One more word for telling, explaining and describing at B1.
\end{grammarlens}

\subsection*{Your turn}
\begin{itemize}
\raggedright
  \item \emph{Say it:} \emph{obruō, obruere}
  \item \emph{Say it:} \emph{obruō, obruere}, once more
  \item \emph{Say it:} a sentence of your own with \emph{obruō, obruere}
  \item \emph{Recall:} say the Latin for shipwrecked, then the Latin for to sink, then say \emph{obruō, obruere} again
\end{itemize}

\begin{tcolorbox}[breakable,colback=teal!4,colframe=teal!35!black,title={Before you move on}]
What does obruō, obruere mean? (\textquotedblleft{}To overwhelm, to bury\textquotedblright{}.) Say it once more.
\end{tcolorbox}
\section[calamitās, calamitātis]{calamitās, calamitātis — a disaster, a misfortune}

\label{lesson:LA-C315-calamitas}



\begin{quote}
Before the new one: say the Latin for to sink, then the Latin for to overwhelm.
\end{quote}

\subsection*{You'll want to know: calamitās, calamitātis}
\textbf{calamitās, calamitātis} — \textquotedblleft{}a disaster, a misfortune\textquotedblright{}.

\begin{quote}
Post illam calamitātem multī cīvēs pauperēs erant.
\end{quote}

\begin{quote}
\emph{After that disaster many citizens were poor.}
\end{quote}

Feminine, 3rd declension. First harm to crops, then any disaster. English \emph{calamity}.

\begin{cousinweb}
\textbf{calamitās} first meant \textquotedblleft{}loss, damage\textquotedblright{}, especially to crops (Lewis \& Short); its origin is not settled. English \emph{calamity} comes from it through French.
\end{cousinweb}

\begin{grammarlens}[title={What you've built}]
One more word for telling, explaining and describing at B1.
\end{grammarlens}

\subsection*{Your turn}
\begin{itemize}
\raggedright
  \item \emph{Say it:} \emph{calamitās, calamitātis}
  \item \emph{Say it:} \emph{calamitās, calamitātis}, once more
  \item \emph{Say it:} a sentence of your own with \emph{calamitās, calamitātis}
  \item \emph{Recall:} say the Latin for to sink, then the Latin for to overwhelm, then say \emph{calamitās, calamitātis} again
\end{itemize}

\begin{tcolorbox}[breakable,colback=teal!4,colframe=teal!35!black,title={Before you move on}]
What does calamitās, calamitātis mean? (\textquotedblleft{}A disaster, a misfortune\textquotedblright{}.) Say it once more.
\end{tcolorbox}
